%% notes-tex/019-simb-multimodal-expression/sections/6-checkpoint.tex
%% SOURCE: experiments/019-simb-multimodal/results/joint_checkpoint_readout.json and the
%% four tables it writes (joint_checkpoint_readout.py, run 2026-10-04 against W&B projects
%% torchcell_019_prot_v19 and torchcell_019_prot_v20); results/
%% expression_morphology_covariation.json and proteome_morphology_covariation.json;
%% conf/cgt_expr_v19_joint_clean.yaml and conf/cgt_expr_v20_conditioned.yaml; the
%% experiment note experiments.019-simb-multimodal.expression-strand-retrospective
%% (entry of 2026-10-04) for job ids and canaries.

\section{Checkpoint, 2026-10-04: the joint round read out, the conditioned round in flight, and the single-deletion triangle\secstatus{todo}}
\label{sec:checkpoint}

\pfstep{The joint round does not support the claim} v19 of \cref{sec:joint-plan} ran as
registered: the both-label store, every head unmasked, three arms from one configuration,
1,200 epochs, one initialization seed per split seed. Eleven of twelve partitions have all
three arms finished (IGB jobs 2413418 on \file{cabbi} and 2413419 on \file{gpu}); partition
11 was still training when this was written. Neither registered test rejects
(\cref{tab:v19-registered}). On the expression head the joint arm is level with the
single-head arm, and the paired difference swings between $-0.045$ and $+0.044$ across
partitions (\cref{tab:v19-partitions}). On the proteome head the joint arm is lower in ten of
eleven partitions on the registered window, past the non-inferiority margin at the point
estimate. One more partition cannot change either reading.
\src{experiments/019-simb-multimodal/results/joint_checkpoint_readout.json}

%% GENERATED by experiments/019-simb-multimodal/scripts/joint_checkpoint_readout.py. Do not edit by hand.
%% SOURCE: W&B zhao-group/torchcell_019_prot_v19 run histories
\begin{table}[htbp]
  \centering
  \footnotesize
  \caption[The v19 registered tests]{The deconfounded joint round on the 11 partitions whose three arms all finished 1,200 epochs, one initialization seed each, paired over split seeds. \emph{single} is the single-head arm (\texttt{K\_expr} or \texttt{K\_prot}) and \emph{joint} is \texttt{K\_joint}; scores are validation \file{pearson_per_feature}. The first two rows are the pre-registered intersection-union test, one-sided at 0.05 each: neither rejects. The $p$ of H1b tests the difference plus the margin. The remaining rows are descriptive: the rolling maximum is the maximum of a centered five-epoch mean, an upward-biased order statistic, and the last two rows read each run at the epoch where its own smoothed validation loss is lowest. \src{experiments/019-simb-multimodal/scripts/joint_checkpoint_readout.py}}
  \label{tab:v19-registered}
  \begin{tabular}{@{}lrrrrrrr@{}}
    \toprule
    score & single & joint & joint $-$ single & SE & positive & $p$ ($t$) & $p$ (sign flip) \\
    \midrule
    H1a expression superiority, epochs 1,000 to 1,200 & 0.094 & 0.093 & $-0.001$ & 0.009 & 5/11 & 0.55 & 0.55 \\
    H1b proteome non-inferiority at 0.015, epochs 200 to 400 & 0.062 & 0.044 & $-0.017$ & 0.006 & 1/11 & 0.64 & 0.64 \\
    proteome, rolling maximum (descriptive) & 0.089 & 0.091 & $+0.002$ & 0.006 & 5/11 & 0.38 & 0.40 \\
    expression, rolling maximum (descriptive) & 0.104 & 0.102 & $-0.002$ & 0.007 & 5/11 & 0.62 & 0.62 \\
    proteome at the validation-loss minimum (descriptive) & 0.037 & 0.004 & $-0.032$ & 0.014 & 3/11 & 0.98 & 0.98 \\
    expression at the validation-loss minimum (descriptive) & 0.050 & 0.031 & $-0.020$ & 0.010 & 4/11 & 0.96 & 0.96 \\
    \bottomrule
  \end{tabular}
\end{table}


%% GENERATED by experiments/019-simb-multimodal/scripts/joint_checkpoint_readout.py. Do not edit by hand.
%% SOURCE: W&B zhao-group/torchcell_019_prot_v19 run histories
\begin{table}[htbp]
  \centering
  \footnotesize
  \caption[v19 by partition]{The registered window scores of the joint round by split seed. Expression is the mean over epochs 1,000 to 1,200 and proteome over 200 to 400. The last three columns give the epoch of the proteome head's rolling maximum in the single-head and the joint arm, and the epoch of the joint arm's smoothed validation-loss minimum: the joint proteome head peaks late, and the loss turns inside the first 110 epochs of a 1,200-epoch run. \src{experiments/019-simb-multimodal/scripts/joint_checkpoint_readout.py}}
  \label{tab:v19-partitions}
  \begin{tabular}{@{}rrrrrrrrrr@{}}
    \toprule
    split & \texttt{K\_expr} & joint & diff & \texttt{K\_prot} & joint & diff & peak, single & peak, joint & loss min \\
    \midrule
    0 & 0.117 & 0.096 & $-0.021$ & 0.122 & 0.073 & $-0.049$ & 276 & 1197 & 80 \\
    1 & 0.069 & 0.108 & $+0.039$ & 0.026 & -0.001 & $-0.027$ & 417 & 988 & 69 \\
    2 & 0.099 & 0.095 & $-0.003$ & 0.117 & 0.073 & $-0.043$ & 341 & 1028 & 109 \\
    3 & 0.121 & 0.092 & $-0.029$ & 0.065 & 0.057 & $-0.007$ & 250 & 450 & 66 \\
    4 & 0.103 & 0.130 & $+0.027$ & 0.073 & 0.050 & $-0.023$ & 549 & 1193 & 87 \\
    5 & 0.080 & 0.083 & $+0.002$ & 0.055 & 0.043 & $-0.012$ & 1177 & 1090 & 74 \\
    6 & 0.096 & 0.115 & $+0.019$ & 0.067 & 0.035 & $-0.031$ & 927 & 244 & 60 \\
    7 & 0.100 & 0.086 & $-0.014$ & 0.038 & 0.054 & $+0.016$ & 151 & 224 & 61 \\
    8 & 0.099 & 0.067 & $-0.032$ & 0.020 & 0.019 & $-0.001$ & 1166 & 1012 & 74 \\
    9 & 0.076 & 0.031 & $-0.045$ & 0.069 & 0.067 & $-0.002$ & 257 & 812 & 30 \\
    10 & 0.071 & 0.115 & $+0.044$ & 0.027 & 0.018 & $-0.010$ & 1190 & 1126 & 55 \\
    \bottomrule
  \end{tabular}
\end{table}


\pfstep{The proteome deficit is at least partly a delay} By the rolling maximum over the full
budget the two arms are level on the proteome head, and the joint arm's head reaches its
maximum late: past epoch 800 in eight of eleven partitions, against four of eleven for the
single-head arm (\cref{tab:v19-partitions}). The window of epochs 200 to 400 was calibrated
on single-head runs in v14 and v16, so it scores an arm that peaks later as an arm that is
worse. The rolling maximum is an upward-biased order statistic, so this reads as no visible
gap at the peak, not as a score. Hypothesis (untested): the expression loss dominates early
optimization of the shared trunk and delays the proteome head.

\pfstep{The loss and the metric disagree in this round as in the long-budget rounds} Every
joint run's smoothed validation loss is lowest between epochs 30 and 109, and the registered
scores are read at epochs 200 to 400 and 1,000 to 1,200 (\cref{tab:v19-partitions}). Read at
each run's own loss minimum, both heads score below their window value in both arms and
the joint arm is lower on both (\cref{tab:v19-registered}, last two rows). This is the
pinball-against-Pearson disagreement of \cref{sec:findings-lossmin} at a shorter budget.
Which stopping rule defines the model is the open decision of \cref{sec:next}, and the
conditioned round is read under both.

\pfstep{The conditioned round, submitted without the masked objective} v20 of
\cref{sec:joint-plan} was changed in one respect before launch: no reveal schedule and no
masked loss. The trainer option \file{multitask.condition_head} reveals one head's measured
labels in full as model input on every step of every stage and drops that head from the loss
and the metrics; \file{condition_permute} reveals another strain's labels, a cyclic shift
inside the batch. Four arms on v19's store, partitions and budget
(\file{conf/cgt_expr_v20_conditioned.yaml}): \texttt{C\_expr} predicts expression with the
proteome revealed, \texttt{C\_prot} predicts the proteome with expression revealed, and
\texttt{C\_exprperm} and \texttt{C\_protperm} are the permuted controls. The unconditioned
controls are v19's \texttt{K\_expr} and \texttt{K\_prot}, already run. IGB job 2423190,
\file{cabbi}, twelve tasks of four runs capped at two concurrent tasks, source commit
\texttt{f09091e2f}. \Cref{tab:v20-partial} is the state at generation and is not a result.

%% GENERATED by experiments/019-simb-multimodal/scripts/joint_checkpoint_readout.py. Do not edit by hand.
%% SOURCE: W&B zhao-group/torchcell_019_prot_v20 and zhao-group/torchcell_019_prot_v19 run histories
\begin{table}[htbp]
  \centering
  \footnotesize
  \caption[v20 while in flight]{PARTIAL: the conditioned round at the epoch each run had reached when this table was generated, of a 1,200-epoch budget. \emph{conditioned} is the mean of the predicted head's validation \file{pearson_per_feature} over the run's latest 20 epochs; \emph{control} is the v19 single-head arm of the same partition over the same epochs; the last column is that control at its registered window. \texttt{C\_expr} predicts expression with the strain's proteome revealed and \texttt{C\_prot} the reverse; the \texttt{perm} arms reveal another strain's measurement. Nothing here is a result until the runs finish. \src{experiments/019-simb-multimodal/scripts/joint_checkpoint_readout.py}}
  \label{tab:v20-partial}
  \begin{tabular}{@{}lrlrrrrr@{}}
    \toprule
    arm & split & predicted & epoch & conditioned & control & diff & control, window \\
    \midrule
    \texttt{C\_expr} & 0 & expression & 1200 & 0.125 & 0.117 & $+0.008$ & 0.117 \\
    \texttt{C\_expr} & 1 & expression & 1200 & 0.127 & 0.083 & $+0.043$ & 0.069 \\
    \texttt{C\_expr} & 2 & expression & 1200 & 0.098 & 0.102 & $-0.004$ & 0.099 \\
    \texttt{C\_expr} & 3 & expression & 1200 & 0.157 & 0.118 & $+0.038$ & 0.121 \\
    \texttt{C\_expr} & 4 & expression & 1199 & 0.152 & 0.114 & $+0.038$ & 0.103 \\
    \texttt{C\_expr} & 5 & expression & 1199 & 0.103 & 0.082 & $+0.022$ & 0.080 \\
    \texttt{C\_expr} & 6 & expression & 299 & 0.048 & 0.035 & $+0.014$ & 0.096 \\
    \texttt{C\_expr} & 7 & expression & 249 & 0.026 & 0.011 & $+0.016$ & 0.100 \\
    \texttt{C\_exprperm} & 0 & expression & 1200 & 0.122 & 0.117 & $+0.005$ & 0.117 \\
    \texttt{C\_exprperm} & 1 & expression & 1200 & 0.003 & 0.083 & $-0.080$ & 0.069 \\
    \texttt{C\_exprperm} & 2 & expression & 1200 & 0.090 & 0.102 & $-0.012$ & 0.099 \\
    \texttt{C\_exprperm} & 3 & expression & 1200 & 0.129 & 0.118 & $+0.011$ & 0.121 \\
    \texttt{C\_exprperm} & 4 & expression & 1199 & 0.085 & 0.114 & $-0.030$ & 0.103 \\
    \texttt{C\_exprperm} & 5 & expression & 1199 & 0.072 & 0.082 & $-0.010$ & 0.080 \\
    \texttt{C\_exprperm} & 6 & expression & 299 & 0.018 & 0.035 & $-0.017$ & 0.096 \\
    \texttt{C\_exprperm} & 7 & expression & 249 & -0.003 & 0.011 & $-0.014$ & 0.100 \\
    \texttt{C\_prot} & 0 & proteome & 1200 & 0.168 & 0.112 & $+0.056$ & 0.122 \\
    \texttt{C\_prot} & 1 & proteome & 1200 & 0.095 & 0.016 & $+0.079$ & 0.026 \\
    \texttt{C\_prot} & 2 & proteome & 1200 & 0.139 & 0.112 & $+0.027$ & 0.117 \\
    \texttt{C\_prot} & 3 & proteome & 1200 & 0.125 & 0.071 & $+0.054$ & 0.065 \\
    \texttt{C\_prot} & 4 & proteome & 1199 & 0.113 & 0.090 & $+0.022$ & 0.073 \\
    \texttt{C\_prot} & 5 & proteome & 1199 & 0.180 & 0.120 & $+0.060$ & 0.055 \\
    \texttt{C\_prot} & 6 & proteome & 299 & 0.077 & 0.069 & $+0.008$ & 0.067 \\
    \texttt{C\_prot} & 7 & proteome & 249 & 0.100 & 0.046 & $+0.054$ & 0.038 \\
    \texttt{C\_protperm} & 0 & proteome & 1200 & 0.126 & 0.112 & $+0.014$ & 0.122 \\
    \texttt{C\_protperm} & 1 & proteome & 1200 & 0.029 & 0.016 & $+0.013$ & 0.026 \\
    \texttt{C\_protperm} & 2 & proteome & 1200 & 0.105 & 0.112 & $-0.007$ & 0.117 \\
    \texttt{C\_protperm} & 3 & proteome & 1200 & 0.088 & 0.071 & $+0.017$ & 0.065 \\
    \texttt{C\_protperm} & 4 & proteome & 1199 & 0.083 & 0.090 & $-0.007$ & 0.073 \\
    \texttt{C\_protperm} & 5 & proteome & 1199 & 0.083 & 0.120 & $-0.037$ & 0.055 \\
    \texttt{C\_protperm} & 6 & proteome & 299 & 0.051 & 0.069 & $-0.018$ & 0.067 \\
    \texttt{C\_protperm} & 7 & proteome & 249 & 0.045 & 0.046 & $-0.001$ & 0.038 \\
    \bottomrule
  \end{tabular}
\end{table}


The conditioned model answers a different question from the joint model. It needs the
conditioning measurement for any strain it is applied to, so it is an imputation claim about
a multimodal model and has to be labeled as one. The reference it has to beat is the ridge of
\cref{tab:triangle}, not the unconditioned arm.

\pfstep{The single-deletion triangle} Each pair of the three single-deletion panels was read
the same way: out-of-fold ridge from one measured modality to the other on the deletions
both measured (\cref{tab:triangle}). The expression-against-morphology side was added on
2026-10-04 with the estimator of the proteome-against-morphology side. Every side carries
signal well above its permuted null, and each is asymmetric in the same direction: the
gene-level panels predict the 278 CalMorph features better than the features predict the
genes. On the morphology sides about half the signal is one shared component: projecting
the first principal component out of both sides takes the moving-feature median from $0.324$
to $0.166$ for expression and from $0.280$ to $0.145$ for the proteome.
\src{experiments/019-simb-multimodal/results/expression_morphology_covariation.json}
\src{experiments/019-simb-multimodal/results/proteome_morphology_covariation.json}

%% GENERATED by experiments/019-simb-multimodal/scripts/joint_checkpoint_readout.py. Do not edit by hand.
%% SOURCE: results/proteome_morphology_covariation.json, results/expression_morphology_covariation.json, review/2026-09-27-joint-review/01_data_ceilings.md
\begin{table}[htbp]
  \centering
  \footnotesize
  \caption[The single-deletion triangle]{Out-of-fold ridge from one measured modality to another on the deletions both panels measured: median held-out Pearson per feature of the predicted modality, five folds over strains. \emph{moving} restricts morphology to the 116 CalMorph features whose replicate reliability is at least 0.5. The proteome and expression rows are the mean per feature from the review of 2026-09-27. The three sides sit on different strain sets and none has been recomputed on a common one. Every strain-permuted null is at most 0.01. \src{experiments/019-simb-multimodal/scripts/joint_checkpoint_readout.py}}
  \label{tab:triangle}
  \begin{tabular}{@{}llrrr@{}}
    \toprule
    given & predicted & all features & moving & shared strains \\
    \midrule
    proteome & expression & 0.226 & -- & 1{,}349 \\
    expression & proteome & 0.218 & -- & 1{,}349 \\
    proteome & morphology & 0.169 & 0.280 & 4{,}314 \\
    morphology & proteome & 0.081 & -- & 4{,}314 \\
    expression & morphology & 0.215 & 0.324 & 1{,}438 \\
    morphology & expression & 0.087 & -- & 1{,}438 \\
    \bottomrule
  \end{tabular}
\end{table}


Against these, the genotype alone reaches $0.101$ on expression and $0.056$ on the proteome
by ridge on the deleted gene's ProtT5 embedding, and adding the genotype to the observed
proteome buys $+0.009$ on expression
(\file{review/2026-09-27-joint-review/01_data_ceilings.md}). The shared structure between
the panels is real and it is in the measurements, not in what a genotype-only trunk sees.

\pfstep{Two structural facts about the model these rounds use} Both bear on morphology more
than on the gene-level heads.
\begin{itemize}
\item \textbf{The CLS token is the wild-type token on this branch.} The encoder runs once on
  the unperturbed graph and the CLS is sliced off before the perturbation operator, so it is
  identical for every strain. Sending the CLS through the operator as one more query
  (\file{perturb_cls}) exists on the gene-interaction line of development (commit
  \texttt{9e2f59551}, 2026-09-09, experiments 025 and 030), where the across-strain standard
  deviation of the perturbed CLS runs $0.31$ to $0.43$. That commit is not an ancestor of
  this branch, and no expression, proteome or morphology run has used it. It makes a global
  readout strain-specific; it does not change the gene tokens, which come out bit-identical.
\item \textbf{The single-key degeneracy is addressed by the null sink, which has one unreplicated run at budget.} With one deleted gene the cross-attention has one key and its
  weights are identically 1. The null sink adds a second key with a learned logit, so
  the weight depends on the querying gene and the query and key parameters receive gradient
  (\file{verify_null_sink_and_pooling.py}, contract C4). Its one trained comparison is v8
  wave 3: four seed pairs stopped at 186 to 189 epochs with the gate bias still near its
  initial $-4$, $+0.0024 \pm 0.0090$, which
  \cref{sec:strand-heads} classifies with every other arm stopped before the validation
  curve leaves its dip. Correction, 2026-10-05: it was also run with the gate open at
  initialization, for 300 epochs on four seed pairs ($+0.0024$, sd $0.0090$) and once for
  4{,}184 epochs (W\&B \texttt{fqk5mbr3}), $0.200$ against $0.210$ for the paired reference
  by rolling maximum. One seed does not resolve a difference of $0.010$.
\end{itemize}

\pfstep{What is open} Whether Figure 3 may rest on a conditioning result, or must be
genotype-only, is the author's decision and sets where the next compute goes. For the
conditioning reading the strongest sides are expression to morphology and proteome to
morphology on the moving features; a store joining the proteome and morphology does not
exist yet, and the observed-label channel places values on gene tokens, so morphology can be
a target and not yet an input. For the genotype-only reading the untried arm is morphology
on its own 4,718 strains read from a perturbed CLS, which changes the training set fourfold
and the readout at once, so the two would have to be separated by arms. Hypothesis
(untested): expression compendia that are not deletion-keyed strengthen the gene-gene
structure the conditioned reads use and leave the genotype-only score where it is.
